% !TEX program = xelatex
\documentclass[12pt,fontset=fandol]{ctexart}
\usepackage[margin=5mm]{geometry}
\usepackage{tikz,amsmath,amssymb,graphicx}
\usepackage[active,tightpage]{preview}
\PreviewEnvironment{tikzpicture}
\setlength\PreviewBorder{5pt}
\usetikzlibrary{arrows.meta,calc}
\pagestyle{empty}
\tikzset{arr/.style={-{Stealth[length=2mm]},semithick}}
\begin{document}
\begin{center}

\begin{tikzpicture}[xscale=1.05,yscale=1]
\draw[arr] (0,0)--(7.4,0) node[right] {$\lambda$};
\draw[arr] (0,0)--(0,3.4) node[above] {相对谱强度};
\draw[thick,domain=.2:7,samples=140]
plot(\x,{.2+2.25*exp(-(\x-2)*(\x-2)/.06)+2.7*exp(-(\x-5)*(\x-5)/.45)});
\draw[dashed] (2,0)--(2,2.5); \draw[dashed] (5,0)--(5,3);
\node[below] at (2,0) {$\lambda_0$}; \node[below] at (5,0) {$\lambda_1$};
\node at (2,3.0) {近弹性峰}; \node at (5,3.5) {康普顿峰};
\node at (3.5,-.9) {峰位与峰宽的定性示意，不作定量读数};
\end{tikzpicture}

\end{center}
\end{document}
